\documentclass[../../../include/open-logic-section]{subfiles}

\begin{document}

\olfileid[hi]{sth}{story}{urelements}
\olsection{मूलतत्त्व हों या नहीं?}

अगले कुछ अध्यायों में हम समुच्चय की संचयी-पुनरावर्ती अवधारणा से अभिगृहीत
निकालने का प्रयास करेंगे। किंतु आगे बढ़ने से पहले हमें \emph{मूलतत्त्वों}
(urelements) के विषय में कुछ और कहना होगा।

\olref[approach]{sec} के चित्र में हम पुराने \emph{समुच्चयों} से ही नए
समुच्चय बना सकते थे। फिर भी हम यह अनुमति देना चाह सकते हैं कि कुछ
\emph{गैर-समुच्चय}---गायें, सूअर, रेत के कण अथवा कुछ और---समुच्चयों के
!!{element}s हो सकते हैं। उस स्थिति में हम कुछ मूल अवयवों, अर्थात
\emph{मूलतत्त्वों}, से आरंभ करते और फिर कहते कि प्रत्येक चरण $S$ पर हम
मूलतत्त्वों अथवा $S$ से पहले के चरणों में बने समुच्चयों (किसी भी संयोजन में)
से बने ``सभी संभव'' समुच्चय बनाएँगे। परिणामी चित्र कुछ ऐसा दिखेगा:
\begin{center}
	\begin{tikzpicture}[scale=0.6]
	\tikzset{cut_here/.style={densely dotted}}
	\draw (-4,6) -- (-1, 0) -- (1,0) -- (4,6); 
	\draw[cut_here] (-5,8)--(-4,6);
	\draw[cut_here] (5,8)--(4,6);
	\draw(-1.5, 1)--(1.5, 1);
	\draw(-2, 2)--(2, 2);
	\draw(-2.5, 3)--(2.5,3);
	\draw(-3, 4)--(3, 4);
	\draw(-3.5, 5)--(3.5, 5);
	\draw(-4, 6)--(4, 6);
	\node[label] at (2, 0) {\small 0};
	\node[label] at (2.5, 1) {\small 1};
	\node[label] at (3, 2) {\small 2};
	\node[label] at (3.5, 3) {\small 3};
	\node[label] at (4, 4) {\small 4};
	\node[label] at (4.5, 5) {\small 5};
	\node[label] at (5, 6) {\small 6};
	\end{tikzpicture}
\end{center}
अब हमें निर्णय करना है: \emph{क्या मूलतत्त्वों की अनुमति देनी चाहिए?}

दार्शनिक दृष्टि से अपने सिद्धांत-निर्माण में मूलतत्त्वों को सम्मिलित करना
सार्थक है। इसका मुख्य कारण अपने समुच्चय सिद्धांत को \emph{अनुप्रयोज्य}
बनाना है। बात स्पष्ट करने के लिए \olref[sfr][siz][]{chap} से याद करें कि
हम दो समुच्चयों $A$ और~$B$ का आकार समान, अर्थात $\cardeq{A}{B}$, तभी और
केवल तभी कहते हैं जब उनके बीच कोई द्विएकी हो। अब यदि मैदान की गायें और बाड़े
के सूअर दोनों समुच्चय बनाते हैं, तो ``गायों की संख्या सूअरों जितनी है'' दावे
का समुच्चय-सैद्धांतिक विवेचन दिया जा सकता है। किंतु यदि हम मूलतत्त्वों को
निषिद्ध कर दें, जिससे गायें और सूअर समुच्चय \emph{नहीं} बनाते, तो वह विवेचन
उपलब्ध नहीं होगा। वास्तव में तब समुच्चय सिद्धांत को स्वयं समुच्चयों के
अतिरिक्त किसी चीज़ पर लागू करने की कोई सीधी क्षमता हमारे पास नहीं होगी।
(मूलतत्त्वों को सम्मिलित करने के और कारणों के लिए
\citealt[pp.~vi, 24, 50--1]{Potter2004} देखें।)

गणितीय व्यवहार में, किंतु, मूलतत्त्वों की अनुमति विरल है। इसका आंशिक कारण
यह है कि मूलतत्त्वों के बिना समुच्चय सिद्धांत सूत्रित करना \emph{बहुत थोड़ा}
सरल है। पर कभी-कभी समुच्चय सिद्धांत से मूलतत्त्वों को बाहर रखने के अधिक
रोचक औचित्य मिलते हैं:
\begin{quote}
	इस विश्वास के अनुरूप कि समुच्चय सिद्धांत गणित की नींव है, हमें केवल
	समुच्चयों की बात करके संपूर्ण गणित को ग्रहण कर सकना चाहिए; अतः हमारे चर
	गाय और सूअर जैसी वस्तुओं पर परासित नहीं होने चाहिए।
	%But if $C$ is a cow, $\{C\}$ is a set, but not a legitimate mathematical object. 
	\citep[p.~8]{Kunen1980}
\end{quote}
इसलिए अनुप्रयोज्यता पर बल मूलतत्त्वों को \emph{सम्मिलित} करने का सुझाव
देगा; अपचायक आधारभूत लक्ष्य (गणित को शुद्ध समुच्चय सिद्धांत में अपचयित करना)
पर बल उन्हें \emph{बाहर रखने} का सुझाव दे सकता है। थोड़ी-सी आलस्यप्रियता भी
मूलतत्त्वों को बाहर रखने की ओर संकेत करती है।

हम सबसे आलस्यपूर्ण मार्ग अपनाएँगे। हालाँकि इसका कुछ शैक्षणिक औचित्य भी है।
हमारा उद्देश्य आपको समुच्चय सिद्धांत के उन तत्त्वों से परिचित कराना है जिनकी
समुच्चय सिद्धांत के दर्शन पर काम आरंभ करने के लिए आवश्यकता होगी। और उस
दार्शनिक साहित्य का अधिकांश भाग मूलतत्त्वों के \emph{बिना} सूत्रित समुच्चय
सिद्धांतों की चर्चा करता है। अतः यदि यह पुस्तक उस साहित्य के काफ़ी निकट रहे,
तो शायद अधिक उपयोगी होगी।

\end{document}
